\Subhead{6.5.1}{Основные задачи нелинейной фильтрации и синтеза систем}

Нелинейность возникает из динамики сообщения, наблюдательной функции или негауссовской статистики. Примеры — слежение за фазой модулированного сигнала, оценка амплитуды при квадратичном детектировании, переключающиеся состояния. Среднее и ковариация уже не обязательно определяют всю апостериорную плотность; несколько гипотез о фазе могут соответствовать одному среднему.

Для марковской модели $\dd x=a(x,t)\dd t+b(x,t)\dd W$ и наблюдения $\dd y=h(x,t)\dd t+\sqrt r\dd V$, независимых шумов и $r>0$ оптимальная фильтрация описывается уравнением Кушнера–Стратоновича:
\begin{equation}
 \dd p=\mathcal L^*p\dd t+
 \frac{h(x,t)-\hat h(t)}r p\,[\dd y-\hat h(t)\dd t],\quad
 \hat h=\int h(x,t)p(x,t)\dd x,
\end{equation}
где $\mathcal L^*p=-\partial_x(ap)+\frac12\partial_x^2(b^2p)$ — оператор Фоккера–Планка. Прогноз распространяет плотность по динамике; инновационный член перевзвешивает состояния согласно наблюдению. Интеграл правой части равен нулю, поэтому нормировка сохраняется при корректных граничных условиях. Это бесконечномерное стохастическое уравнение, а не простая замена матриц Калмана функциями.

Для дискретных наблюдений логика особенно ясна:
\begin{align}
 p(x_k|y_{1:k-1})&=\int p(x_k|x_{k-1})p(x_{k-1}|y_{1:k-1})\dd x_{k-1},\\
 p(x_k|y_{1:k})&\propto p(y_k|x_k)p(x_k|y_{1:k-1}).
\end{align}
После нормировки выбирают оценку по цене ошибки: среднее для квадратичной, максимум плотности для соответствующей иной задачи или решение о гипотезе.

Приближённые реализации: расширенный Калман линеаризует около оценки и пригоден при достаточно узкой почти гауссовской плотности; сигма-точки приближают перенос моментов; частицы представляют плотность взвешенными реализациями. У сильной нелинейности, многомодальности и срывов фазового слежения локальная гауссовская аппроксимация может терять нужную гипотезу.

Синтез приёмника включает выбор состояния, модели канала/шума, достаточной статистики, критерия и реализуемой рекурсии. Проверяют не только локальную дисперсию, но и вероятность грубых ошибок, ложного захвата, устойчивость к ошибке модели и вычислительную задержку. Численный алгоритм не становится оптимальным только от названия «байесовский».
